% Current experimental section generated from frozen run artifacts on 2026-09-01.
% Requires \usepackage{amsmath,booktabs,graphicx}.
% The citation keys below are intentionally generic; map them to the paper's .bib.

\section{Experiments}
\label{sec:experiments}

We evaluate Nash Rejection Sampling (Nash-RS) in a common online-preference
learning pipeline and compare it with method-native implementations of six
baselines.  This section reports every completed comparison available at the
time of writing.  We distinguish the small-scale implementation and stability
gate from the current main experiment, and we explicitly mark incomplete
results.  In particular, the Qwen2.5-3B COMAL run is still queued and is not
included in any 3B comparison.

\subsection{Experimental setup}
\label{sec:experimental-setup}

\paragraph{Data.}
We use prompts from the training split of \textsc{UltraFeedback}.  We normalize
whitespace, remove exact duplicates by SHA-256 of the normalized prompt, and
retain prompts whose rendered chat template contains at most 128 tokens.  A
deterministic shuffle with seed 20260827 produces three disjoint files: 2,048
training prompts, 256 selection prompts, and 1,024 test prompts.  The 0.5B
stability gate uses the first 512 prompts of the frozen 2,048-prompt training
file, whereas the 3B experiment uses all 2,048 training prompts.  Neither the
selection nor test prompts are used for optimization.  Unless stated
otherwise, responses are limited to 512 new tokens.

\paragraph{Models and optimization.}
The stability gate starts from \texttt{Qwen/Qwen2.5-0.5B-Instruct}; the main
experiment starts from \texttt{Qwen/Qwen2.5-3B-Instruct} at revision
\texttt{aa8e72537993ba99e69dfaafa59ed015b17504d1}.  All reported runs update all
model parameters; none uses LoRA or another parameter-efficient adapter.  We
use a policy learning rate of $10^{-6}$, except for MPO, for which the
method-native implementation uses $5\times10^{-7}$.  OpenRLHF PPO runs use a
critic learning rate of $10^{-5}$.  The learning-rate schedule is constant and
has no warm-up.  Every method receives the same prompt pool, base model,
response-length limit, preference oracle, and seed within a comparison.  The
number of optimizer steps differs because we preserve each method's native
update schedule: OpenRLHF PPO consumes two rollout prompts per global step,
TRL Nash-MD and EGPO consume one prompt per step, and MPO performs two inner
steps per outer prompt.

\paragraph{General-preference oracle.}
Training and evaluation use an equal-weight mixture of two independently
trained scalar preference models: (i)
\texttt{OpenAssistant/reward-model-deberta-v3-large-v2}, revision
\texttt{c355404e}, with effective temperature 2.0609; and (ii)
\texttt{Skywork/Skywork-Reward-V2-Qwen3-0.6B}, revision
\texttt{d35d68a7}, with effective temperature 3.9374.  The first model scores a
pair-encoded input and the second scores a chat-formatted input.  On a
32-prompt diagnostic set with four responses per prompt, the two components
disagree on 30.73\% of pairs.  The aggregate preference has no observed
three-cycle on this small diagnostic set, but is not exactly representable by
a single Bradley--Terry model: its logit and probability root-mean-square
residuals are 0.00910 and 0.00202, respectively.  Consequently, the benchmark
tests general-preference optimization even though the observed aggregate is
close to Bradley--Terry on this diagnostic sample.

\paragraph{Compared methods and implementation fidelity.}
Table~\ref{tab:method-fidelity} records the exact entry used for each method.
This is our \emph{Main} comparison---Nash-RS against practical,
method-native baselines---and not the separate controlled OpenRLHF/PPO
ablation.  Reward-PPO and Nash-RS intentionally use OpenRLHF because PPO is
their declared practical optimizer.  The remaining methods retain their
native update rules.

\begin{table}[t]
  \centering
  \caption{Method-native implementations used in the current experiments.}
  \label{tab:method-fidelity}
  \small
  \resizebox{\linewidth}{!}{%
  \begin{tabular}{lll}
    \toprule
    Method & Backend & Update implemented in our experiments \\
    \midrule
    Reward-PPO & OpenRLHF PPO & Scalar mixture reward followed by PPO \\
    Self-play & TRL 0.13 NashMDTrainer & Nash-MD endpoint with geometric-mixture coefficient zero \\
    Nash-MD & TRL 0.13 NashMDTrainer & Token-level geometric-mixture Nash-MD-PG \\
    Nash-RS & OpenRLHF PPO & Gibbs rejection sampling, implicit reward, then PPO \\
    MPO & Safe-RLHF-based native loop & MPO Algorithm~1, ReMax, proximal and periodic magnet updates \\
    EGPO & Authors' trainer & Online-IPO extragradient prediction/correction \\
    COMAL & Authors' pipeline & Sampling, scoring, log-probability pipeline, and INPO inner updates \\
    \bottomrule
  \end{tabular}}
\end{table}

For Nash-RS we use $B_1=B_2=2$, proposal batches of four, and the correct
acceptance event
\begin{equation}
  u \leq \exp\!\left(-\widehat g/\tau\right)
  \quad\Longleftrightarrow\quad
  \widehat g \leq -\tau\log u,
\end{equation}
where $u\sim\mathrm{Unif}[0,1]$.  The 0.5B three-seed gate uses $	au=1$;
the 3B experiment uses $	au=0.5$, selected by the ablation in
Section~\ref{sec:tau-ablation}.

\paragraph{Evaluation.}
We generate one greedy response from every policy on each of the 256 frozen
selection prompts.  For methods $i$ and $j$, let
$W_{ij}$ be their mean pairwise preference probability over prompts.  We report
\begin{align}
  \mathrm{AvgWin}(i) &= \frac{1}{M-1}\sum_{j\ne i} W_{ij}, \\
  \mathrm{WorstWin}(i) &= \min_{j\ne i} W_{ij}, \\
  \widehat{\mathrm{Exploit}}(i) &=
    \max\{0,\tfrac12-\mathrm{WorstWin}(i)\}.
\end{align}
Thus, lower empirical exploitability is better.  Confidence intervals resample
prompts with 10,000 bootstrap replicates and fixed bootstrap seed 20260829.
The evaluation oracle is the same frozen preference mixture used for training;
we discuss this limitation in Section~\ref{sec:exp-limitations}.

\subsection{0.5B three-seed implementation and stability gate}
\label{sec:0p5b-gate}

Table~\ref{tab:0p5b-main} aggregates seeds 47, 101, and 211.  Each policy sees
512 training prompts and each seed is evaluated on the same 256 selection
prompts.  Reward-PPO has the highest mean average and worst-case win rates in
this gate.  Nash-RS with the untuned value $\tau=1$ does not outperform the
baselines on average.  This result motivated the temperature ablation below.
Because the experiment was registered as an implementation/stability gate,
we do not treat it as the final paper comparison.

\begin{table}[t]
  \centering
  \caption{Qwen2.5-0.5B stability gate over three seeds.  Entries are
  mean $\pm$ sample standard deviation across seeds.  Higher is better for
  win rates; lower is better for exploitability.}
  \label{tab:0p5b-main}
  \small
  \begin{tabular}{lccc}
    \toprule
    Method & Average win rate & Worst-case win rate & Exploitability \\
    \midrule
    Base model & $0.4968\pm0.0096$ & $0.4753\pm0.0316$ & $0.0247\pm0.0316$ \\
    Reward-PPO & $\mathbf{0.5208\pm0.0319}$ & $\mathbf{0.5122\pm0.0289}$ & $\mathbf{0.0042\pm0.0073}$ \\
    Self-play & $0.5020\pm0.0086$ & $0.4800\pm0.0228$ & $0.0205\pm0.0222$ \\
    Nash-MD & $0.4958\pm0.0038$ & $0.4743\pm0.0268$ & $0.0257\pm0.0268$ \\
    Nash-RS ($\tau=1$) & $0.4957\pm0.0086$ & $0.4738\pm0.0309$ & $0.0262\pm0.0309$ \\
    MPO & $0.4901\pm0.0033$ & $0.4698\pm0.0260$ & $0.0302\pm0.0260$ \\
    EGPO & $0.4993\pm0.0012$ & $0.4774\pm0.0249$ & $0.0226\pm0.0249$ \\
    COMAL & $0.4995\pm0.0097$ & $0.4775\pm0.0318$ & $0.0225\pm0.0318$ \\
    \bottomrule
  \end{tabular}
\end{table}

The per-seed variability is material.  For example, Reward-PPO average win
rates are 0.5117, 0.4945, and 0.5562 for seeds 47, 101, and 211, respectively;
the corresponding Nash-RS values are 0.5022, 0.4989, and 0.4860.  The base
model itself varies across generation seeds, so the across-seed summary is
more informative than any single 0.5B run.

\paragraph{Training cost.}
Table~\ref{tab:0p5b-cost} reports allocated GPU time rather than the internal
sample-timer field: all jobs use one H200, so GPU-hours equal PBS walltime in
hours.  PM calls count individual component-model calls; because the oracle has
two components, one logical pairwise query generally contributes two PM calls.
Nash-RS uses approximately 4.1 times as many PM calls as Nash-MD and 2.4 times
as many generated tokens, reflecting the real cost of rejection sampling with
$B_1=B_2=2$.  Nevertheless, its end-to-end walltime is lower than Nash-MD in
this implementation because OpenRLHF/vLLM batches generation more efficiently
than the TRL method-native backend.

\begin{table}[t]
  \centering
  \caption{Qwen2.5-0.5B training cost over 512 prompts.  Values are means over
  three seeds; GPU-hours are allocated H200 walltime.}
  \label{tab:0p5b-cost}
  \small
  \begin{tabular}{lrrrr}
    \toprule
    Method & Optimizer steps & PM calls & Generated tokens (k) & H200-hours \\
    \midrule
    Reward-PPO & 256 & 512 & 145.0 & 0.402 \\
    Self-play & 512 & 1,024 & 313.3 & 2.382 \\
    Nash-MD & 512 & 1,024 & 312.3 & 2.363 \\
    Nash-RS & 256 & 4,193 & 753.7 & 1.366 \\
    MPO & 1,024 & 2,048 & 476.9 & 1.708 \\
    EGPO & 512 & 1,024 & 320.3 & 0.798 \\
    COMAL & 512 & 1,024 & 320.8 & 1.203 \\
    \bottomrule
  \end{tabular}
\end{table}

Each 0.5B evaluation makes 14,336 component PM calls.  The three evaluation
jobs require 0.110, 0.134, and 0.126 H200-hours, respectively; evaluation cost
is not included in Table~\ref{tab:0p5b-cost}.

\subsection{Rejection-temperature ablation}
\label{sec:tau-ablation}

We retrain Nash-RS at
$\tau\in\{0.25,0.5,0.75,1,2,4\}$ with seed 47 and the same 512-prompt and
256-step budget.  Table~\ref{tab:tau-sweep} shows that acceptance and compute
are strongly related, whereas policy quality is non-monotone in $\tau$.  The
acceptance rate rises from 17.9\% at $\tau=0.25$ to 89.3\% at $\tau=4$, but
$\tau=0.5$ has the best average win rate, worst-case win rate, and empirical
exploitability.  We therefore use $\tau=0.5$ at 3B.  The low-temperature
$\tau=0.25$ run uses a maximum-proposal cap of 1,024 rather than 64 so that it
can complete the same accepted-sample budget; this change affects cost but not
the target acceptance rule.

\begin{table}[t]
  \centering
  \caption{Nash-RS temperature ablation at 0.5B (seed 47).  All rows use 512
  prompts and 256 optimizer steps.}
  \label{tab:tau-sweep}
  \small
  \resizebox{\linewidth}{!}{%
  \begin{tabular}{rrrrrrrr}
    \toprule
    $\tau$ & Acceptance & PM calls & Tokens (k) & H200-hours & Avg. win & Worst win & Exploit. \\
    \midrule
    0.25 & 0.179 & 12,492 & 1,513.7 & 2.851 & 0.4884 & 0.4644 & 0.0356 \\
    0.50 & 0.415 & 5,964 & 943.5 & 1.736 & \textbf{0.5271} & \textbf{0.5122} & \textbf{0.0000} \\
    0.75 & 0.578 & 4,568 & 790.6 & 1.421 & 0.5137 & 0.4878 & 0.0122 \\
    1.00 & 0.665 & 4,102 & 762.3 & 1.384 & 0.4936 & 0.4702 & 0.0298 \\
    2.00 & 0.808 & 3,560 & 720.3 & 1.309 & 0.5118 & 0.4863 & 0.0137 \\
    4.00 & 0.893 & 3,318 & 729.8 & 1.330 & 0.5134 & 0.4872 & 0.0128 \\
    \bottomrule
  \end{tabular}}
\end{table}

The 95\% prompt-bootstrap interval for the average win rate of $\tau=0.5$ is
$[0.5138,0.5398]$.  Corresponding intervals for $\tau=0.75$, 2, and 4 are
$[0.5001,0.5272]$, $[0.4988,0.5244]$, and $[0.5016,0.5254]$.  These intervals
quantify prompt uncertainty for one training seed and must not be interpreted
as across-training-seed confidence intervals.

\subsection{Qwen2.5-3B main comparison}
\label{sec:3b-main}

We next train full-parameter Qwen2.5-3B policies on all 2,048 frozen training
prompts using seed 47.  Table~\ref{tab:3b-main} reports the six completed
methods.  Nash-RS obtains the highest average win rate (0.5286), the only
worst-case win rate above 0.5 (0.5154), and zero empirical exploitability
within the evaluated policy set.  Its bootstrap interval excludes 0.5 for both
average and worst-case win rate.  Reward-PPO is second by average win rate, but
its worst-case win rate remains below 0.5.  Because this experiment currently
has only one training seed and lacks COMAL, we regard it as strong preliminary
evidence rather than a final multi-seed result.

\begin{table}[t]
  \centering
  \caption{Current Qwen2.5-3B main result on 256 selection prompts (seed 47).
  Brackets are 95\% prompt-bootstrap intervals.}
  \label{tab:3b-main}
  \small
  \resizebox{\linewidth}{!}{%
  \begin{tabular}{lccc}
    \toprule
    Method & Average win rate & Worst-case win rate & Exploitability \\
    \midrule
    Base & $0.4947\ [0.4878,0.5014]$ & $0.4710\ [0.4594,0.4824]$ & $0.0290\ [0.0176,0.0406]$ \\
    Reward-PPO & $0.5102\ [0.4994,0.5212]$ & $0.4846\ [0.4726,0.4965]$ & $0.0154\ [0.0035,0.0274]$ \\
    Self-play & $0.4951\ [0.4884,0.5020]$ & $0.4713\ [0.4597,0.4828]$ & $0.0287\ [0.0172,0.0403]$ \\
    Nash-MD & $0.4847\ [0.4771,0.4920]$ & $0.4623\ [0.4506,0.4738]$ & $0.0377\ [0.0262,0.0494]$ \\
    Nash-RS & $\mathbf{0.5286\ [0.5189,0.5383]}$ & $\mathbf{0.5154\ [0.5035,0.5269]}$ & $\mathbf{0.0000\ [0.0000,0.0000]}$ \\
    MPO & $0.4927\ [0.4850,0.5004]$ & $0.4693\ [0.4568,0.4819]$ & $0.0307\ [0.0181,0.0432]$ \\
    EGPO & $0.4939\ [0.4870,0.5007]$ & $0.4699\ [0.4586,0.4812]$ & $0.0301\ [0.0188,0.0414]$ \\
    \bottomrule
  \end{tabular}}
\end{table}

Table~\ref{tab:3b-pairwise} gives the complete pairwise matrix.  Nash-RS beats
every other row-wise opponent in expectation, including Reward-PPO (0.5154),
Nash-MD (0.5377), MPO (0.5307), and EGPO (0.5301).

\begin{table}[t]
  \centering
  \caption{Qwen2.5-3B pairwise win-probability matrix $W_{ij}$; rows are the
  evaluated policy and columns are the opponent.}
  \label{tab:3b-pairwise}
  \scriptsize
  \resizebox{\linewidth}{!}{%
  \begin{tabular}{lrrrrrrr}
    \toprule
    & Base & R-PPO & Self & N-MD & N-RS & MPO & EGPO \\
    \midrule
    Base       & .5000 & .4869 & .4996 & .5086 & .4710 & .5017 & .5005 \\
    Reward-PPO & .5131 & .5000 & .5130 & .5216 & .4846 & .5149 & .5141 \\
    Self-play  & .5004 & .4870 & .5000 & .5087 & .4713 & .5022 & .5011 \\
    Nash-MD    & .4914 & .4784 & .4913 & .5000 & .4623 & .4933 & .4918 \\
    Nash-RS    & .5290 & .5154 & .5287 & .5377 & .5000 & .5307 & .5301 \\
    MPO        & .4983 & .4851 & .4978 & .5067 & .4693 & .5000 & .4988 \\
    EGPO       & .4995 & .4859 & .4989 & .5082 & .4699 & .5012 & .5000 \\
    \bottomrule
  \end{tabular}}
\end{table}

The two component rewards agree with the pairwise result.  Nash-RS has mean
OpenAssistant-DeBERTa score 1.9600 and mean Skywork-Qwen score 4.0889, compared
with 1.8928 and 3.7165 for Reward-PPO.  The base-model scores are 1.6835 and
3.6149.  The 95\% intervals are $[1.7011,2.2193]$ and $[3.7089,4.4777]$ for
Nash-RS, and $[1.6344,2.1537]$ and $[3.3189,4.1108]$ for Reward-PPO.

\paragraph{Training and evaluation cost.}
Table~\ref{tab:3b-cost} reports actual one-H200 PBS walltime.  Nash-RS requires
6.34 times as many logical component PM calls and 2.84 times as many generated
tokens as Nash-MD.  Despite this sampling overhead, Nash-RS completes in 11.87
H200-hours versus 13.62 for Nash-MD, again because the backends have different
generation throughput.  This is a practical end-to-end comparison, not an
isolated algorithmic complexity comparison.  The common seven-model 3B
evaluation takes 0.104 H200-hours, 5,376 logical pairwise comparisons, and
10,752 component PM calls.

\begin{table}[t]
  \centering
  \caption{Qwen2.5-3B training cost on 2,048 prompts (seed 47).  All runs use
  one H200; H200-hours therefore equal PBS walltime.}
  \label{tab:3b-cost}
  \small
  \begin{tabular}{lrrrr}
    \toprule
    Method & Optimizer steps & PM calls & Generated tokens (M) & H200-hours \\
    \midrule
    Reward-PPO & 1,024 & 2,048 & 0.761 & 4.453 \\
    Self-play & 2,048 & 4,096 & 1.426 & 13.528 \\
    Nash-MD & 2,048 & 4,096 & 1.429 & 13.620 \\
    Nash-RS & 1,024 & 25,972 & 4.052 & 11.871 \\
    MPO & 4,096 & 8,192 & 2.104 & 10.292 \\
    EGPO & 2,048 & 4,096 & 1.429 & 6.272 \\
    \bottomrule
  \end{tabular}
\end{table}

Average generated response lengths at evaluation are 346.0 tokens for the
base model, 372.9 for Reward-PPO, 346.2 for self-play, 346.1 for Nash-MD, 357.1
for Nash-RS, 343.3 for MPO, and 346.2 for EGPO.  Corresponding truncation rates
are 43.8\%, 45.3\%, 43.4\%, 43.8\%, 42.6\%, 42.6\%, and 42.2\%.  These high
rates motivate a longer-generation sensitivity check in future experiments.

\subsection{Development studies}
\label{sec:development-studies}

Before freezing the method-native comparison, we ran three checks that shaped
the final protocol.

\paragraph{Short multi-seed Nash-RS pilot.}
A 24-step 0.5B pilot evaluated the final checkpoint against the base model on
128 independent prompts for seeds 47, 101, and 211.  Direct win rates are
0.5063, 0.5075, and 0.4978, giving $0.5039\pm0.0053$ across seeds.  A
10,000-sample hierarchical bootstrap gives the 95\% interval
$[0.4822,0.5250]$.  The mean scalar-reward improvement is
$0.0195\pm0.0331$, with interval $[-0.0917,0.1280]$.  This pilot therefore
verified end-to-end learning but did not establish a significant improvement.

\paragraph{UltraFeedback 2K checkpoint selection.}
In an OpenRLHF development run, checkpoint selection on 256 prompts chose
global step 400.  On the disjoint 1,024-prompt test set, this checkpoint beats
the base model with win probability 0.5378 and 95\% prompt-bootstrap interval
$[0.5296,0.5462]$; its empirical exploitability within the two-policy set is
zero.  The two mean component rewards change from 0.3936 and $-0.5720$ for the
base model to 0.8009 and 0.0641 after training.  This experiment used the same
preference models for optimization and evaluation and is therefore a pipeline
validation rather than independent external validation.

\paragraph{Longer general-preference pilot.}
A 96-step learning-curve experiment on 16 validation prompts selects step 72
by average win rate (0.5400).  On 128 independent prompts, however, average win
rates are 0.4903, 0.4984, and 0.5113 for the base, step-72, and step-96 models,
respectively; the corresponding 95\% intervals overlap.  This discrepancy is
why the main protocol freezes a larger 256-prompt selection set and avoids
claiming monotone improvement from small learning curves.

\subsection{Limitations of the current evidence}
\label{sec:exp-limitations}

First, the 0.5B three-seed experiment is an implementation gate, not a final
benchmark.  Second, the 3B comparison currently contains only one training
seed, and COMAL is pending because its two-H200 job has not entered the queue;
we do not impute or extrapolate its result.  Third, all reported win rates use
the same frozen two-component preference mixture used for training.  This
controls evaluation across methods but may reward oracle-specific
over-optimization; independent reward models and human or public benchmark
judges are required for the final paper.  Fourth, 3B truncation rates exceed
42\%, so length sensitivity remains a relevant confound.  Finally, prompt
bootstrap intervals quantify uncertainty over evaluation prompts conditional
on a trained policy; they do not replace multi-seed uncertainty over training.
For these reasons, the main quantitative conclusion supported by the current
artifacts is deliberately narrow: at Qwen2.5-3B, seed 47, and the frozen
2,048/256 train/selection protocol, Nash-RS is the strongest of the six
completed methods under the shared general-preference oracle, while its
rejection sampler incurs substantially more PM calls and generated tokens.

